\documentclass{article}
\usepackage{enumitem}
\usepackage[margin=1in]{geometry}
\usepackage{natbib}
\usepackage{graphicx} % Required for inserting images

\title{From Data to Dynamics: Discovering Governing Equations for Ocean–Climate Variability}
\author{Anand Gnanadesikan, Somdatta Goswami, Jan Drgona}
\date{January 2026}

\begin{document}

\maketitle

Many key oceanographic and climate phenomena involve emergent behavior spanning a range of scales and processes. We argue that recent advances in machine learning, particularly in physics-informed learning and equation discovery to create a timely opportunity to revisit these longstanding problems from a new perspective. The core idea is to move beyond black-box emulation towards learning underlying equations themselves: compact, interpretable dynamical models that can be analyzed, compared, and embedded within existing Earth system frameworks. We do so by bringing together three research groups to work on three use cases for which there is evidence that such equation sets can be found, but where there is no general agreement on what they should be.

The three problems are the following:
\begin{enumerate}[leftmargin=*]
    \item 
    Global distribution of phytoplankton biomass. Sitting at the base of the marine food web, phytoplankton control how much nutrients and energy are available for global fisheries. However, many current generations of Earth System Models are unable to simulate the distribution of key functional groups, limiting our ability to predict potential impacts on fisheries due to climate variability and change. However, recent work in the Gnanadesikan group at JHU \citep{dutta2025using,gnanadesikan2025machine,holder2023well,jin2024using} suggests that these distributions are, in fact, highly predictable from a relatively small set of environmental conditions and that it is the relationships between the environmental conditions and target biological variables that are inaccurate. However, the relationships isolated to date have involved complex machine learning models rather than more readily interpretable equations. 
    \item 
    Variability in high latitude convection, particularly in the Atlantic, is known to exert a strong impact on fisheries and is believed to affect hurricane activities, floods, and droughts throughout North America and Europe. Here, too, however, Earth System Models vary greatly in the amplitude and natural period of variability simulated. While a number of simplified models have been proposed to explain the time-mean structure of the overturning and its variability, there is no overall theory that allows us to understand what sets the amplitude and period of variability. Theories that we have found to work in the Southern Ocean \citep{gnanadesikan2020feedbacks} do not appear to work in the North Atlantic.
    \item The El Ni\~no-Southern Oscillation system is the dominant mode of interannual variability on the planet and has major impacts on floods, drought, and fire worldwide. As with the previous two phenomena, this is also an area where models disagree in terms of amplitude, period, and degree of phase-locking. While it is quite easy to find linear delay-oscillator equations for ENSO \citep{russell2014understanding,vanOldenburgh2005searching}, it is difficult to find {\it robust} equations that can be used across models and observations to understand differences and develop predictions. 
\end{enumerate}

We propose to attack these three problems with new ML methods currently being developed in the Civil and Systems Engineering Group at JHU to discover equations in complex systems. We have picked the three systems above since PI Gnanadesikan has worked on developing simple models for each of them. While he has found that there is enough information in the system to build simple models that explain certain aspects of each phenonemon, he is also very aware of the shortcomings of such models. Additionally, all three systems have important societal impacts.

The first of the ML research groups, Jan Drgona's group at JHU, currently leads development of the  Neuromancer SciML library \citep{drgona2023domain} in PyTorch for solving constrained optimization, physics-informed machine learning, and optimal control problems. The group has used these methods to understand control of heating and cooling in buildings, where one seeks to model the response of a system containing a potentially large number of reservoirs exchanging with each other at unknown rates\citep{das2024machine,drgona2020all}. This is in many ways analogous to the challenge of constructing dynamical box models of the coupled system, where it is unclear which centers of action in the ocean and atmosphere group together. Additionally, this package can be used for equation discovery, and work by Gnanadesikan's graduate student has begun to apply it to ocean biogeochemical modeling. Preliminary work suggests that neural operators can skillfully simulate the evolution of observed phytoplankton biomass, and we aim to build on this work to replace the neural operator with closed-form equations.

The second group, Somdatta Goswami's group at Johns Hopkins, has focused on equation discovery for canonical dynamical systems using neural-operator frameworks that decouple functional-form identification from parameter estimation. Building on this foundation, we plan to employ a Deep Operator Network (DeepONet) architecture to uncover compact, interpretable dynamical equations linking atmospheric circulation patterns to ocean overturning and gyre variability in the North Atlantic. DeepONet provides a flexible framework for learning mappings between high-dimensional spatio-temporal fields, such as atmospheric forcing or sea-surface temperature anomalies, and corresponding dynamical responses or modal amplitudes. Because these datasets are inherently high-dimensional and multi-scale, we will first apply reduced-order modeling techniques, e.g., an autoencoder, to identify low-dimensional latent spaces that capture the dominant structures across different spatial and temporal scales. Within these latent manifolds, we will deploy a deep hidden-physics model to infer distinct sets of governing equations characterizing dynamics at each scale. By integrating physics-aware regularization with sparsity-promoting inversion, the learned operator enables (i) hidden-physics discovery, via symbolic or sparse regression on the surrogate to extract governing terms, and (ii) parameter identification, through efficient inverse estimation of closure coefficients across model resolutions \citep{goswami2025121284,kag2024learning,mandl2025separable}. The resulting framework will yield compact parametric representations with quantified uncertainties that can be systematically compared across resolutions and validated for dynamical fidelity before reintegration into a coupled climate model. 

We believe that this approach will yield a number of important benefits. First, insofar as we can find families of interpretable governing equations this will greatly increase our mechanistic understanding of these emergent processes. Second, such equations can also allow cross-model comparison of dynamical regimes, either because the underlying equations appear to be different or the constants within those equations are different. Third, the biogeochemical modeling task may serve to provide reduced-order equations for ocean biogeochemistry within Earth System Models, as a template for developing such subcomponents for other aspects of the Earth System. Finally, it seems likely that a better understanding of the underlying dynamics will lead to better ability to predict variability and change of phenomena that matter to society at large.

\bibliographystyle{plain}
\bibliography{project.bib}
\end{document}
